\documentclass{article}
\usepackage{../assignment_preamble}

\title{Homework 9}
\author{Ravi Kini}
\date{December 12, 2025}

\begin{document}

\maketitle

\problem{}
At the end of the Cretaceous period, a mass extinction event wiped out most of the plant and animal species on Earth, including all non-avian dinosaurs. One proposed explanation was the eruption of the Deccan traps, a large shield volcano. The other explanation, which is currently widely-accepted by the paleontological community, is the impact of a massive asteroid. Evidence includes the presence of an impact crater in the Yucatán Peninsula that exhibits disproportionately high concentrations of iridium. The iridium has also been dated with high accuracy to the same time period as the extinction event.

\clearpage
\problem{}
\subproblem{(a)}
The holdout method evaluates the performance of a model by splitting the data set into a training set and a testing set. The training set is used to develop the model, and the testing set is used to measure how well the model adapts to unseen data.
\subproblem{(b)}
Both LOOCV and $n$-fold CV involve partitioning the data set, and rotating the partition usd as the testing set. The primary difference is that LOOCV has a single data point as the testing set, while $n$-fold CV has one of $n$ equal partitions as the testing set. LOOCV can be treated as a special case of $n$-fold CV, where $n$ is the size of the data set.

\clearpage
\problem{}
For gradient descent to be applied to a parametric model $M$, it must be possible to change the values of the parameters continuously. This is not applicable to default logic, which is necessarily discrete; neither the set of default rules nor the ordering over the rules can be changed continuously. Further, the measure of lack of fit does not change continuously, even assuming the parameters can be changed continuously, as the set of sentences that are skeptical non-monotonic consequences of the premise are discrete as well.
\end{document}